% ================= TÓM TẮT =================
\newpage
\thispagestyle{plain}

\begin{center}
    {\bfseries\Large TÓM TẮT}
\end{center}

\vspace{1cm}

Phát hiện biển báo giao thông là một khối chức năng bắt buộc của hệ thống hỗ trợ lái tiên tiến và
xe tự hành. Bài toán đặt ra hai ràng buộc đối nghịch: mô hình phải đủ chính xác để không bỏ sót
biển báo, đồng thời phải đủ nhẹ và đủ nhanh để chạy thời gian thực trên phần cứng nhúng gắn trên
xe. Trong vài năm gần đây, hai họ kiến trúc thống trị lĩnh vực phát hiện đối tượng đi theo hai
triết lý khác hẳn nhau: họ một giai đoạn dựa trên mạng tích chập mà đại diện là YOLO, và họ
transformer dự đoán theo tập hợp mà đại diện là DETR. Câu hỏi thực tiễn đặt ra cho người triển
khai là: với một bộ dữ liệu quy mô vừa và ngân sách tính toán của một đồ án học phần, họ nào cho
kết quả tốt hơn, và cái giá phải trả cho mỗi lựa chọn là gì.

Đề tài sử dụng bộ dữ liệu \textit{pkdarabi/cardetection} trên Kaggle, gồm 4.969 ảnh độ phân giải
416$\times$416 với 6.012 hộp bao thuộc 15 lớp là đèn tín hiệu và biển giới hạn tốc độ. Toàn bộ quy
trình được xây dựng dưới dạng các mô-đun Python độc lập, có kiểm thử tự động, với sổ tay Jupyter
chỉ đóng vai trò điều phối. Bốn nhóm thực nghiệm được tiến hành: so sánh mô hình nền giữa YOLOv8n
và DETR-ResNet50; khảo sát ba kiến trúc nhẹ thay thế gồm YOLO11n, SSDLite-MobileNetV3 và
D-FINE-Nano; lượng tử hoá sau huấn luyện theo ba đường FP16, ONNX INT8 và OpenVINO INT8; và cuối
cùng là chưng cất tri thức từ YOLO26s sang YOLO26n.

Mọi mô hình được huấn luyện với cùng một hạt giống ngẫu nhiên và được đánh giá trên tập kiểm tra
tách biệt. Các độ đo gồm mAP tại ngưỡng IoU 0,5 và mAP trung bình trên dải ngưỡng 0,5--0,95, kèm
theo độ trễ suy luận, tốc độ khung hình và dung lượng mô hình để phản ánh ràng buộc triển khai.
Một phân tích khám phá dữ liệu đầy đủ được thực hiện trước khi huấn luyện, bao gồm phân bố lớp,
thống kê kích thước hộp bao, bản đồ nhiệt vị trí vật thể và kiểm tra chất lượng nhãn.

Kết quả cho thấy YOLOv8n đạt mAP@0,5 bằng 0,9539 và mAP@0,5:0,95 bằng 0,8058 trên tập kiểm tra,
với 6,25~MB trọng số và độ trễ 8,83~ms mỗi ảnh. DETR-ResNet50 chỉ đạt mAP@0,5 bằng 0,1220 sau 10
chu kỳ huấn luyện, tức kém hơn gần tám lần, đồng thời nặng hơn 26,6 lần và chậm hơn 7,2
lần. Ba kiến trúc nhẹ ở nhóm thực nghiệm thứ hai đều không vượt được mô hình nền. Lượng tử hoá
FP16 giữ nguyên độ chính xác trong khi giảm khoảng một nửa dung lượng, còn INT8 tuy nén mạnh hơn
nhưng lại chậm đi trên môi trường thử nghiệm. Chưng cất tri thức cho kết quả \emph{âm}: mô hình
học trò được chưng cất đạt 0,9014, thấp hơn chính nó khi huấn luyện thông thường (0,9348).

Đóng góp của đề tài không nằm ở một con số mAP cao mà ở hai phát hiện mang tính phương pháp luận.
Thứ nhất, khoảng cách giữa DETR và YOLO trong điều kiện này chủ yếu do điều kiện huấn luyện bất
đối xứng chứ không phải do bản chất kiến trúc, và báo cáo phân tách rõ từng yếu tố gây nhiễu. Thứ
hai, phân tích thành phần bộ dữ liệu chỉ ra rằng 38,7\% số ảnh huấn luyện là ảnh cắt cận cảnh có
cạnh biển chiếm tới 66,8\% cạnh ảnh, khiến phân bố kích thước vật thể bị lưỡng cực và giải thích
vì sao mô hình đạt mAP rất cao trên tập kiểm tra nhưng suy giảm rõ rệt trên video đường phố thực
tế. Hạn chế lớn nhất là hai nhánh mô hình dùng hai thư viện tính mAP khác nhau và bộ dữ liệu tồn
tại hiện tượng trùng ảnh giữa các tập.

\vspace{0.5cm}
\noindent\textbf{Từ khoá:} phát hiện biển báo giao thông, YOLO, DETR, lượng tử hoá,
chưng cất tri thức, dịch chuyển miền, vật thể nhỏ.
